\documentclass{article}
\usepackage{amsmath,amssymb}
\begin{document}
\section*{The simultaneous rare/balanced corner of the half-tax inequality}
Put \(s=4q(1-q)\), \(c=4q^2\), and
\[
R=M+\tfrac12\Delta-J,
\quad J=h_2(sp)-p h_2(s),
\]
\[
M=(1-sp)h_2\!\left(\frac{cp(1-p)}{1-sp}\right)+sp h_2(p),
\quad \Delta=p\log_2\frac{1-p}{p}.
\]
The frozen candidate on this two-bit gate is \(R\ge0\).

Let \(h(x)=-x\ln x-(1-x)\ln(1-x)\) and
\(\mathcal R=(\ln2)R\). Uniformly for \(q\) in every compact
subinterval of \((0,1)\),
\[
\mathcal R(p,q)=p\left[(c+\tfrac12-s)\ln\frac1p+C(q)\right]
+O\!\left(p^2\ln\frac1p\right),                                      \tag{1}
\]
where
\[
C(q)=c(1-\ln c)-s(1-\ln s)+h(s)
=c-s-c\ln c-(1-s)\ln(1-s).                                           \tag{2}
\]
Indeed \(h(ap)=ap(\ln(1/p)-\ln a+1)+O(p^2)\) uniformly here, and
\(cp(1-p)/(1-sp)=cp+O(p^2)\). Substitution gives (1), retaining all
normalization terms.

The logarithmic coefficient factors exactly:
\[
c+\tfrac12-s=8(q-\tfrac14)^2,                                        \tag{3}
\]
while
\[
C(1/4)=\ln2-\tfrac12>0.                                               \tag{4}
\]
Therefore, for every sequence \(p_k\downarrow0\), \(q_k\to1/4\), and
\(\eta_k=q_k-1/4\),
\[
\frac{R(p_k,q_k)}{p_k}
=8\eta_k^2\log_2\frac1{p_k}
+1-\frac1{2\ln2}+o(1).                                                \tag{5}
\]
Consequently
\[
\liminf_{k\to\infty}\frac{R(p_k,q_k)}{p_k}
\ge1-\frac1{2\ln2}=0.2786524795\ldots>0.                              \tag{6}
\]
Thus an open collar \(0<p<\varepsilon\), \(|q-1/4|<\delta\) satisfies
\(R>0\). This is a local two-bit theorem; it does not prove the remaining
bias strips, higher arity, nonreusing transport, universal GU, or FEI.
\end{document}
